\section{A proposed research programme}
\label{sec:questions}

The results suggest several precise continuations. The questions below are
not assertions that the expected extension must be true. Each includes an
initial route and a criterion for what would count as progress.

\subsection{Uniform resonance and Stieltjes finite parts}
\begin{enumerate}[label=R\arabic*.,leftmargin=*]
\item \textbf{Complex approach to the singular argument.}
Extend Theorem~\ref{res:main} from positive $t$ to complex sectors, with
an explicit branch of $\log t$ and a region avoiding the nearest singularities
$\pm2\pi i$. The pole-pairing identity still provides a candidate normal
form. The substantive issue is a common majorant and a useful condition
on the complex centered scale. A first deliverable is a sector theorem
with constants uniform on closed subsectors.

\item \textbf{The shift boundary $a\downarrow0$.}
Determine a uniform form when $a$, $t$, and the spectral displacement
approach their singular values together. Subtract the elementary $n=0$
term before expanding $G_a$: generalized Stieltjes constants themselves
become singular at that boundary. The goal is to identify the additional
dimensionless ratio and prove an expansion that matches both the compact
$a$ theorem and the isolated first summand.

\item \textbf{Growing spectral jets.}
The present estimates allow every fixed number of order derivatives.
Find the maximal growth range of the derivative index for which a
quantitative finite-part expansion remains effective. Cauchy bounds on
the centered Gamma quotient provide factorial estimates, but optimizing
the contour radius against the derivative order and $T$ is necessary.
This would turn the fixed-jet Stieltjes identities into a joint
large-index asymptotic calculus.

\item \textbf{Growing transition windows.}
Replace the fixed compact set of $u$ in Theorem~\ref{res:main} by a
window increasing with $T$. Specify whether the desired error is absolute
or relative, since $e^{-u}$ behaves differently in opposite half-planes
and the leading profile has nonzero imaginary zeros. The exact expression
\eqref{res:pair} gives a starting point for explicit estimates in terms of
$|u|/T$. A useful first result would identify maximal real windows for
each fixed truncation order and then treat complex sectors separately.
\end{enumerate}

\subsection{Shifted Gamma and polygamma calculus}
\begin{enumerate}[label=G\arabic*.,leftmargin=*]
\item \textbf{Concavity beyond the first Stieltjes primitive.}
For which $r\geq2$ is the squared increment $E_{r,r}(a)$ concave on
$(0,1)$? The all-order formula \eqref{gamma:eq:all-jet-expansion} gives
both the cusp and the analytic coefficients, but the first-level
coefficient sign proof does not automatically persist. A rigorous
counterexample would be as informative as a classification. Compute
symbolic coefficient signs and use certified bounds before proposing
a universal higher-order conjecture.

\item \textbf{A distributional polygamma extension.}
Choose a specific periodic finite-part normalization for $\psi$ and its
argument derivatives, then extend the translated correlation kernel.
Determine the delta-function terms at integral shifts and prove that
argument differentiation commutes with the chosen normalization.
The Fourier formula gives the spectral candidate; the endpoint terms
are the essential missing proof.

\item \textbf{Three-factor correlations and rational shifts.}
Develop integrals of three translated Hurwitz jets. Their Fourier
coefficients obey a relation $n_1+n_2+n_3=0$, leading to constrained
double sums and a natural comparison with the manuscript's Tornheim
layer. At rational shifts, determine explicit character decompositions
and minimal formal constant spaces at fixed weight and jet order.
Any assertion about independence of the resulting values requires a
separate arithmetic argument.

\item \textbf{Uniform certificates for mixed cusps.}
Generalize Corollary~\ref{gamma:cor:tail} from $E_{1,1}$ to $E_{r,m}$,
with a bound explicit in all three variables $a,r,m$. A useful target
is an error bound that remains practical after reducing the shift to
$0\leq a\leq1/2$. Such a result would also support interval evaluation
of the higher concavity questions.
\end{enumerate}

\subsection{The Herglotz transition and its arithmetic profile}
\begin{enumerate}[label=H\arabic*.,leftmargin=*]
\item \textbf{Matching the small-ratio and fixed-argument regimes.}
Combine the exact gamma mixture with the Bessel formula for $S$ to
obtain a remainder estimate that stays useful when $\lambda=x/r$
approaches zero with $r$. The phase of the exponentially small term
changes across the gamma distribution. A compact-ratio Taylor bound
does not capture that phase at arbitrary shrinking ratios. The desired
theorem should recover the existing fixed-$x$ oscillation and the
new transition profile from a single controlled representation.

\item \textbf{Quantitative imaginary-zero geometry.}
Theorem~\ref{herg:thm:profilezeros} gives exactly one simple centered zero
between successive imaginary poles. Determine the location $b_N$
within $(2\pi N,2\pi(N+1))$ as $N\to\infty$, including the effect of
the divisor weights. Prove explicit bounds before seeking a limiting
fractional position: the residues fluctuate arithmetically. A related
goal is a canonical product with justified convergence and normalization.

\item \textbf{Weighted profiles and higher polylogarithm primitives.}
Replace $n^{-2}$ in the positive profile by $n^{-\beta}$ for $\beta>1$.
The same elementary kernel suggests divisor weights of a different
order, a changed Mellin product, and an iterated primitive ladder.
Determine exactly which properties survive: complete ratio expansions,
imaginary-axis interlacing, Bessel representations, and closed normalized
polylogarithm antiderivatives. Treat analytic continuation in $\beta$
separately from the positive-measure range.

\item \textbf{Rigorous numerical evaluation over the full axis.}
Implement ball-arithmetic evaluation that switches between the
large-ratio zeta expansion and the small-ratio Bessel expansion using
proved total error bounds. For the original finite-$r$ function,
combine the central-moment certificate with a direct Hurwitz method.
The stopping rule must account for rounding error as well as omitted
terms, especially when measuring exponentially small oscillations.
\end{enumerate}

\subsection{Formal jets, module structure, and special values}
\begin{enumerate}[label=A\arabic*.,leftmargin=*]
\item \textbf{Three active generators.}
For a local Frobenius algebra, classify ideals for which the homology
profile is $(\ell,3\ell,3\ell,\ell)$ and those for which the middle
dimension is larger or smaller. The explicit example
$(xy,xz,yz)$ shows that colength alone is insufficient. A first
tractable class is monomial ideals in rectangular jet algebras, where
the first syzygies admit a combinatorial presentation.

\item \textbf{Beyond abelian Smith factors.}
Determine the $B$-module extensions in the reflected integral quotients,
including when the two signs have isomorphic torsion modules and when
their torsion-free quotients are free over $B$. The annihilator example
\eqref{ns:eq:module-obstruction} supplies a small test case. The correct
invariants must retain module generators and extension data, not just
dimensions after reduction modulo two.

\item \textbf{Arithmetic realizability of deformations.}
Which mixed active sequences arise from actual families of polylogarithm
or Hurwitz spectral jets after choosing an integral coefficient lattice?
Construct the specialization map explicitly and track which primes are
inverted. This separates the abstract Frobenius-algebra theorem from
questions about values of analytic functions.

\item \textbf{The unresolved $S_6$ special value.}
Retain the manuscript's exact definition and current conjecture as the
target; the results above do not decide it. A productive next step is
to combine the existing integral reductions with a precisely specified
space of multiple-polylogarithm periods, prove the functional reductions
in that space, and only then test any proposed basis relation.
High-precision integer-relation evidence may guide a candidate identity,
but cannot supply the missing reduction or independence proof.
\end{enumerate}

The first integration priorities are the two resolved extensions, the
strict convexity theorem with its computable series, and the explicit
three-generator obstruction. The next analytic priority is a uniform
matching theorem for the Herglotz endpoints; the next algebraic priority
is a syzygy description beyond two generators. These are concrete
continuations of the proofs given here.

